\addtocounter{chapter}{-1}
\chapter{Motivation}\label{parMotivation}%
\addcontentsline{itc}{chapter}{\protect\numberline{\thechapter}Motivation}

\section{Some results on Ising's model}\label{parSomeResultsOnIsingsModel}

In this subsection we recall the definition of Ising's model
and give two classical results on it, namely Theorems~\ref{c3793} and~\ref{c6177}.
In \S~\ref{parProblematics}, considerations on these results
will serve as a motivation to the work of this monograph.


\subsection{Ising's model}\label{parIsingsModel}

Ising's celebrated model is a basic model of equilibrium thermodynamics,
which represents a ferromagnetic material:
%
\begin{Def}
For~$n$ an integer, consider the lattice $\ZZ^n$ endowed
with its usual graph structure (each vertex has $2n$ neighbours),
and denote by~$\dist$ the graph distance.
Define $\Omega = \{\pm1\}^{\ZZ^n}$,
and for~$\vec\omega\in\Omega$, set formally:
\[\label{for9073} H(\vec\omega) = - \frac{1}{2} \sum_{\dist(i,j)=1} \omega_i\omega_j .\]
%
Then, for~$T\geq0$, the \emph{Ising model on~$\ZZ^n$ at temperature $T$}
is, formally, a probability measure~$\Pr$ on~$\Omega$ such that
$\Pr[\vec\omega] \propto \exp\big(-T^{-1}H(\vec\omega)\big)$.
In rigorous terms, saying that $\Pr$ is an equilibrium measure for Ising's model
means that for all~$i\in\ZZ^n$,
for all~$\vechat\omega_{\C{\{i\}}} \in \{\pm1\}^{\C{\{i\}}}$,
\[\label{f1762} \Pr\big[ \omega_i = \hat\omega_i
\big| \vec\omega_{\C{\{i\}}} = \vechat\omega_{\C{\{i\}}} \big]
\propto \exp\Big( T^{-1} \sum_{\dist(i,j)=1} \hat\omega_i \hat\omega_j \Big) .\]
\end{Def}

Ising's model and the phase transition it exhibits
have been the subject of dozens of works; see~\cite{Grimmett} for an overview.
Here we are interested in the subcritical regime:
%
\begin{Thm}[Subcritical regime, \cite{Peierls}]
There is a $T_{\mrm{c}}<\infty$ (the `Curie temperature')
such that the solution of~(\ref{f1762}) is unique for $T>T_{\mrm{c}}$ .
\end{Thm}

For $T>T_{\mrm{c}}$ one says that they are in the \emph{subcritical regime}.
An interesting feature of this regime is that for distant~$i$ and~$j$,
the random variables~$\omega_i$ and~$\omega_j$ are `almost independent'.
That phenomenon, called \emph{exponential decay of correlations},
is stated by the following theorem:
%
\begin{Thm}[Exponential decay of correlations, \cite{sharp_transition}]\label{t3795}
For Ising's model on~$\ZZ^n$ in the subcritical regime,
\begin{ienumerate}
\item\label{i6505} For all~$i\in\ZZ^n$, $\Pr[\omega_i=-1] = \Pr[\omega_i=1] = 1/2$.
\item\label{i1366} There exists~$\psi>0$ and~$C<\infty$
such that for all~$i,j\in\ZZ^n$,
\[ | \EE[\omega_i\omega_j] | \leq C \exp\big(-\psi\dist(i,j)\big) .\]
\end{ienumerate}
\end{Thm}

\subsection{Absence of \texorpdfstring{$\beta$}{beta}-mixing}

Theorem~\ref{t3795} states that two distant spins~$i$ and~$j$ are exponentially decorrelated.
However, it does not inform us about the dependence between `bunches' of spins.
The question is the following: if $I$ and~$J$ are two disjoint, distant subsets of~$\ZZ^n$,
to what extent are~$\vec\omega_I$ and~$\vec\omega_J$ independent?

To answer such a question, the first thing to do is to define a way of measuring independence
between `complicated' variables having an arbitrarily large range
like~$\vec\omega_I$ and~$\vec\omega_J$.
The most common choice is the $\beta$-mixing coefficient:
%
\begin{Def}\label{def3167}\strut
\begin{enumerate}
\item Recall that for~$\mu$, $\nu$ two probability measures
on the same measurable space $(\Omega,\mcal{F})$,
the \emph{total variation distance} between~$\mu$ and~$\nu$ is
the total mass of both the positive and the negative parts of the signed measure~$\nu-\mu$, that is,
$\dist_{\mathrm{TV}}(\mu,\nu) = \sup_{A\subset\mcal{F}} | \nu(A) - \mu(A) |$.
\item If $X$ and~$Y$ are two random variables (with arbitrary ranges)
defined on the same space, then one defines the \emph{$\beta$-mixing coefficient}
between~$X$ and~$Y$ as
\[ \beta(X,Y) \coloneqq \dist_{\mathrm{TV}} ( \Law_X\otimes\Law_Y\,,\,\Law_{(X,Y)} ) .\]
\end{enumerate}
\end{Def}
%
Notice that $\beta(X,Y)$ actually only depends on the $\sigma$-algebras%
~$\sigma(X)$ and~$\sigma(Y)$~\cite[Formula~(1.5)]{Bradley-review}.
The following proposition is immediate:
%
\begin{Pro}\label{pr7428}\strut
\begin{ienumerate}
\item One has always $\beta(X,Y) \in [0,1]$, and
\item $\beta(X,Y) = 0$ if and only if $X$ and~$Y$ are independent ;
\item $\beta(X,Y) = 1$ if and only if $\Law_X\otimes\Law_Y$
and~$\Law_{(X,Y)}$ are mutually disjoint.
\item\label{i7428d} If $X'$ is $X$-measurable and $Y'$ is $Y$-measurable,
then $\beta(X',Y') \leq \beta(X,Y)$.
\end{ienumerate}
\end{Pro}
%
So, $\beta(X,Y)$ is a way of measuring `how much $X$ and~$Y$ are correlated'.

With that tool at hand, decorrelation between bunches of spins
in statistical physics models has already been thoroughly studied.
Concerning Ising's model, there are two well-known great results:
%
\begin{Thm}[Weak mixing property, \cite{Martinelli}]\label{t4774}
For Ising's model on~$\ZZ^n$ in the subcritical regime,
there exists~$\psi>0$ and~$C<\infty$ (the same as in Theorem~\ref{t3795})
such that for all disjoint $I,J\subset\ZZ^n$:
\[\label{f7869}
\beta \big( \vec\omega_I , \vec\omega_J \big) \ \leq
\ C \sum_{(i,j)\in I\times J} \exp\big(-\psi \dist(i,j)\big) .\]
\end{Thm}

\begin{Thm}[Complete analyticity, \cite{complete_analyticity}]\label{t4505}
There exists some $T_{\mrm{c}}\leq T_{\mrm{c}}'<\infty$%
\footnote{\label{Tc'=Tc?}It is not known whether $T_{\mrm{c}}'=T_{\mrm{c}}$ today,
but in general situations weak mixing does not always imply complete analyticity.
A classical counterexample is Ising's model with external field~\cite[\S~2]{MO}.}
such that for $T \ab {>T_{\mrm{c}}'}$,
Ising's model is \emph{completely analytical}, i.e.\ 
there exists~$\psi'>0$ and~$C'>0$ such that the following holds:
for all~$K\subset\ZZ^n$, for all `boundary' conditions $\vechat\omega_K \in \{\pm1\}^K$,
denoting $\Pr_{\vechat\omega_K} \coloneqq \Pr[\Bcdot|\vec\omega_K=\vechat\omega_K]$,
Formula~(\ref{f7869}) holds with~$\Law$ replaced by~$\Law_{\vechat\omega_K}$
and~$\psi, C$ replaced resp.\ by~$\psi'$ and~$C'$.
\end{Thm}

Thanks to Theorem~\ref{t4774}, we get an exponential decay of correlation
between two bunches of spins of fixed size when the distance between these bunches increases.
However, we cannot say much about decorrelation between bunches of \emph{variable} size
which are at fixed distance from each other.
For instance, for $n=2$ fix $x>0$ and define
$I_\ell \coloneqq \{(0,y) \colon |y|\leq \ell\}$, resp.\ $J_\ell \coloneqq \{(x,y) \colon {|y|\leq \ell}\}$.
Then Theorem~\ref{t4774} gives us something like:
\[\label{f4961}
\beta \big( \vec\omega_{I_\ell} , \vec\omega_{J_\ell} \big) \lesssim
\ C \ell e^{-\psi x} .\]
But recall that a $\beta$-mixing coefficient is always bounded by~$1$;
so, for $\ell \gtrsim e^{\psi x}/C$, (\ref{f4961}) tell us absolutely nothing about
the decorrelation between~$I_\ell$ and~$J_\ell$.

Though the bound~(\ref{f7869}) is not completely optimal,
the previous point is an \emph{intrinsic} shortcoming of $\beta$-mixing coefficients,
in the sense that it can be proved that bounds like~(\ref{f4961}) \emph{must} become trivial when
$\ell\longto\infty$:
%
\begin{Thm}\label{c3793}
For all~$T_{\mrm{c}}<T<\infty$, for all~$x>0$,
denoting $I \coloneqq \{0\}\times\ZZ$ and $J \coloneqq \{x\}\times\ZZ$,
one has
\[ \beta \big(\vec\omega_I , \vec\omega_J \big) = 1 .\]
\end{Thm}

\begin{proof}
Denote $i_0 \coloneqq (0,0)$, resp.\ $j_0 \coloneqq (x,0)$.
As we told in Theorem~\ref{t3795}-(\ref{i6505}),
$\EE[\omega_{i_0}], \EE[\omega_{j_0}] = 0$.
Interpretation of Ising's model as a random-cluster model~\cite[\S~1.4]{Grimmett}
shows that $\Pr[\omega_{i_0}=\omega_{j_0}] > 1/2$,
so we define
\[ \gamma \coloneqq \EE[\omega_{i_0}\omega_{j_0}] > 0 .\]

Now, let~$N$ be some large integer, fixed for the time being,
let~$p$ be some large integer and define
$i_1, \ldots, i_N$, resp.~$j_1, \ldots, j_N$,
by~$i_k \coloneqq (0,kp)$, resp.\ $j_k \coloneqq (x,kp)$;
denote by~$P_{N,p}$ the joint law of~%
$(\omega_{i_1},\ldots,\omega_{i_N},\omega_{j_1},\ldots,\omega_{j_N})$.
By translation invariance, for each~$k$,
$(\omega_{i_k},\omega_{j_k})$ has the same law as $(\omega_{i_0},\omega_{j_0})$,
which law we denote by $P_1$.
%
Then when~$p \longto \infty$, by Theorem~\ref{t4774},
$P_{N,p}$ tends to the law $P_{N,\infty} \coloneqq P_1^{\otimes N}$.%
\footnote{Note that $P_{N,p}$
takes its values in a space of finite dimension,
so there is no ambiguity when speaking of its convergence.}
In other words, $P_{N,\infty}$ is the law such that
all the~$(\omega_{i_k},\omega_{j_k})$ are independent
with $P_{N,\infty}[\omega_{i_k}=\eta \enskip\text{and}\enskip \omega_{j_k}=\theta]
= (1+\gamma\eta\theta)/4$ for all~$k$.
%
Therefore, the value of~$\beta\big((\omega_{i_k})_k, (\omega_{j_k})_k\big)$
under the law $P_{N,p}$, which by Proposition~\hbox{\ref{pr7428}-(\ref{i7428d})}
is a lower bound for~$\beta \big(\vec\omega_I , \vec\omega_J \big)$,
tends to its value under~$P_{N,\infty}$ when~$p \longto \infty$.
This is summed up by the following formula:
\[\label{fo7691} \beta \big(\vec\omega_I, \vec\omega_J\big) \geq \beta_{P_{N,\infty}}
\big((\omega_{i_k})_{1\leq k\leq N}, (\omega_{j_k})_{1\leq k\leq N}\big) .\]

To end the proof, we will bound the right-hand side of~(\ref{fo7691}) below
by a quantity which tends to~$1$ when~$N\longto\infty$.
Denote by~$\tilde{P}_{N,\infty}$ the product of two the marginals
of~$P_{N,\infty}$ relative resp.\ to the~$(\omega_{i_k})_k$ and the~$(\omega_{j_k})_k$,
so that $\beta_{P_{N,\infty}} \big((\omega_{i_k})_k, (\omega_{j_k})_k\big)
= \dist_{\mathrm{TV}} (P_{N,\infty},\tilde{P}_{N,\infty})$
by the very definition of the $\beta$-mixing coefficient.
Obviously the expression of~$\tilde{P}_{N,\infty}$ is the same
as the expression of~$\tilde{P}_{N,\infty}$, but with~$\gamma$ replaced by~$0$ in the definition of~$P_{N,\infty}$.
Under $P_{N,\infty}$, $(\omega_{i_k}\omega_{j_k})_{1\leq k\leq N}$
is a sequence of i.i.d.\ random variables having a certain law with mean $\gamma > \gamma/2$,
so that by the law of large numbers,
\[ P_{N,\infty} \bigg[N^{-1}\sum_{k=1}^N \omega_{i_k}\omega_{j_k} \leq \frac{\gamma}{2} \bigg]
\stackrel{N\longto\infty}{\longto} 0 .\]
Similarly, since $0 < \gamma/2$,
\[ P_{N,\infty} \bigg[N^{-1} \sum_{k=1}^N \omega_{i_k}\omega_{j_k} \leq \frac{\gamma}{2} \bigg]
\stackrel{N\longto\infty}{\longto} 1 ,\]
so that
\[ \dist_{\mathrm{TV}} \big(P_{N,\infty},\tilde{P}_{N,\infty}\big) \geq
\bigg| \tilde{P}_{N,\infty}
\bigg[ N^{-1} \sum_{k=1}^N \omega_{i_k}\omega_{j_k} \leq \frac{\gamma}{2} \bigg] - P_{N,\infty}
\big[ \text{\it the same} \big] \bigg|
\stackrel{N\longto\infty}{\longto} 1 ,\]
which proves our point.
\end{proof}



\subsection{Presence of \texorpdfstring{$\rho$}{rho}-mixing}

So, Theorem~\ref{c3793} tells us that, for Ising's model on~$\ZZ^2$,
there is a `full' correlation between~$\vec\omega_{I}$ and~$\vec\omega_{J}$
in the sense of $\beta$-mixing.
Yet it is well known too that
Theorem~\ref{t4774} nevertheless implies a Hilbertian form of decorrelation
(called ``$\rho$-mixing'', cf.\ Remark~\ref{rmk8718}) between these variables:
%
\label{parc6177}\begin{Thm}\label{c6177}
For Ising's model on~$\ZZ^2$ in the subcritical regime,
defining as before $I = \{0\}\times\ZZ$ and $J = \{x\}\times\ZZ$ for some $x>0$,
one has for all~$f\in\ldb(\vec\omega_I)$ and~$g\in\ldb(\vec\omega_J)$:
\[\label{f2881} |\EE[fg]| \leq e^{-\psi x} \ecty(f)\ecty(g) ,\]
where $\psi$ is the same as in Theorem~\ref{t4774}.
\end{Thm}

\begin{proof}
Define the operator
\[ \begin{array}{rrcl} P \colon & \ldb(\vec\omega_I) & \to & \ldb(\vec\omega_J) \\
& f & \mapsto & f^{\sigma(\vec\omega_J)} .\end{array} \]
(Recall that $f^{\sigma(\vec\omega_J)}$ is an alternative notation for~$\EE[f|\vec\omega_J]$,
insisting on the its being a $\sigma(\vec\omega_J)$-measurable function).
Then (\ref{f2881}) is equivalent to proving that
$\VERT P \VERT \leq e^{-\psi x}$ (see \S~\ref{parOperatorInterpretation}).
Now for all~$t \in \{0,\ldots,x\}$, denote $\omega_{(t)} \coloneqq \vec\omega_{\{t\}\times\ZZ}$,
and for all~$t \in \{1,\ldots,x\}$,
\[ \begin{array}{rrcl} \pi_t \colon & \ldb(\omega_{(t-1)}) & \to & \ldb(\omega_{(t)}) \\
& f & \mapsto & f^{\sigma(\omega_{(t)})} .\end{array}\]
Due to the fact that the interactions in Ising's model have only range~$1$,
$\omega_{(0)} \to \omega_{(1)} \to \cdots \to \omega_{(x)}$
is a Markov chain, and therefore
\[\label{f3677} P = \pi_x \circ \cdots \circ \pi_2 \circ \pi_1 .\]
Now, by horizontal translation all the~$\ldb(\omega_{(t)})$ can be identified
with a common Hilbert space $H$.
Then all the~$\pi_t$ are identified with operators on~$\ldb(H)$,
and by the translation invariance of the model all these operators are actually the same.
$P$ is also identified with an operator on~$\ldb(H)$, and (\ref{f3677}) becomes:
\[ P = \pi^x .\]
But $\pi$ is self-adjoint because,
as the model is invariant by translation \emph{and by reflection},
the Markov chain $\omega_{(0)} \to \cdots \to \omega_{(x)}$
is stationary \emph{and reversible}.
In particular $\pi$ is a normal operator, and thus $\VERT P \VERT = \VERT \pi \VERT^x$.
So, proving that $\VERT P \VERT \leq e^{-\psi x}$
is equivalent to proving that $\VERT \pi \VERT \leq e^{-\psi}$,
which will be our new goal.

Take $C<\infty$ like in Theorem~\ref{t4774}.
For~$\ell$ an integer, denote~$I_\ell$ and~$J_\ell$ to be resp.~%
$\{0\}\times\{-\ell,\ldots,\ell\}$ and~$\{x\}\times\{-\ell,\ldots,\ell\}$.
Let~$f$ be a bounded%
\footnote{In fact here it is superfluous to impose that $f$ is bounded since
$\vec\omega_{I_\ell}$ can only take a finite number of values.
I wrote the proof like this just to underline
that the finiteness of the range of the~$\omega_i$ does not play any role in the proof.}
function of~$\ldb(\vec\omega_{I_\ell})$ and denote $M \coloneqq \|f\|_{L^\infty}$.
By translation, $f$ can also be identified with
a function of~$\ldb(\vec\omega_{J_\ell})$,
which is also bounded by~$M$.
%
Now, since
\[ \EE \big[ f(\vec\omega_{I_\ell})f(\vec\omega_{J_\ell}) \big] =
\Cov \big( f(\vec\omega_{I_\ell}) \,,\, f(\vec\omega_{J_\ell}) \big) \\
= \int f(\vec\omega_{I_\ell})f(\vec\omega_{J_\ell}) \,
\dx{\big(\Law(\vec\omega_{I_\ell\uplus J_\ell})} -
\Law(\vec\omega_{I_\ell})\otimes\Law(\vec\omega_{J_\ell})\big) ,
\]
we can apply~(\ref{f7869}) to~$I_{\ell}$ and~$J_{\ell}$ to obtain:
\[\label{f3079} | \EE[f(\vec\omega_{I_\ell})f(\vec\omega_{J_\ell})] |
\leq M^2\cdot 2C(2\ell+1)^2 e^{-\psi x} .\]

In terms of operators, (\ref{f3079}) means that
\[\label{f3079'} \big|\langle f, Pf \rangle_{\ldb(H)}\big| \leq 2(2\ell+1)^2M^2Ce^{-\psi x} .\]
As the value of~$x$ played no particular role to establish~(\ref{f3079'}),
that formula can be generalized
into
\[ \big|\langle f, \pi^t f \rangle_{\ldb(H)}\big| \leq 2(2\ell+1)^2M^2 C e^{-\psi t} \]
for all~$t\in\NN^*$.
Letting $t$ tend to infinity, we obtain that for all~$\ell$,
for all~$f\in\ldb(\vec\omega_{I_\ell}) \cap L^\infty$,
\[ \limsup_{t\longto\infty} \big( \ln |\langle f, \pi^t f \rangle| \big)^{1/t} \leq e^{-\psi} .\]
But $\bigcup_{\ell\in\NN} \big(\ldb(\vec\omega_{I_\ell})\cap L^\infty\big)$ is a dense subset of
$\ldb(\vec\omega_I)$, so by Lemma~\ref{l3253} set in appendix,
we conclude that $\VERT \pi \VERT_{\ldb(H)} \leq e^{-\psi}$, which is what we wanted.
\end{proof}

\begin{Rmk}
Claim~\ref{c3793} and Theorem~\ref{c6177} adapt straighforwardly,
with similar proofs, to any~$n\geq2$,
replacing~$I$ by~$\{0\}\times\ZZ^{n-1}$ and~$J$ by~$\{x\}\times\ZZ^{n-1}$.
\end{Rmk}



\section{Problematics}\label{parProblematics}

Thanks to Theorem~\ref{c6177}, we see that the Hilbertian concept of $\rho$-mixing
can reveal some independence between infinite bunches of lowly correlated variables
in situations where the $\beta$-mixing coefficient does not show any independence at all.
In the proof we gave, $\rho$-mixing appeared as a corollary
of $\beta$-mixing for finite bunches of spins.
What additional hypotheses did we need to get our corollary?
We used at least the following:
%
\begin{itemize}
\item To introduce the Markov chain $\omega_{(0)} \to \cdots \to \omega_{(x)}$,
we used that the interactions of our model had finite range.
\item To identify all the spaces $\ldb(\omega_{(t)})$,
we used that $I$ and~$J$ had the same shape and that one could tile up $\ZZ^2$
with a sequence of tiles having that shape (namely, here, tiles of the form $\{t\}\times\ZZ$).
\item To say that all the~$\pi_t$ were the same modulo that identification,
we used the translation invariance of the model.
\item To state that the stationary Markov chain $\omega_{(0)} \to \cdots \to \omega_{(x)}$
was reversible, we used the reflection invariance of the model.
\item To use Lemma~\ref{l3253}, we used the exponential decay of correlations.
\end{itemize}
%
All these points make the proof of Theorem~\ref{c6177} we gave
in~\S~\ref{parSomeResultsOnIsingsModel} quite difficult to generalize.
What, for instance, if we take~$I$ and~$J$ with arbitrary shapes,
just requiring that $\dist(I,J) \geq x$?
What if we consider statistical physics models with infinite-range interactions? Etc..
The above arguments would not work any more! Yet, we do not have the impression that
the presence of $\rho$-mixing fundamentally relies
on the peculiar symmetries of the case we treated\dots

So, here will be the goal of this work:
\emph{establishing $\rho$-mixing estimates by general methods}.
To achieve this goal, I shall try to concentrate on the properties of $\rho$-mixing
`for itself', rather than on its links with other forms of decorrelation.
I will carry out a thorough study of the $\rho$-mixing coefficient,
in order to get $\rho$-mixing results for `complicated' variables
from decorrelation results \emph{of the same type} for more `basic' variables;
in other words, I will \emph{tensorize} Hilbertian decorrelations.
It turns out that tensorization for such correlation coefficients gives results which are
quite robust as the size of the bunches of variables increases.
Thanks to this method, I shall obtain fairly new decorrelation theorems
for various models of statistical physics.

This monograph is intended to be complete in some sense.
I mean, besides the core of this work—namely, tensorization results—,
I have tried to answer several other questions which appeared natural to me
concerning Hilbertian decorrelation.
This includes studying many examples,
finding sharp criteria for maximal decorrelation,
looking at the optimality issues in the tensorization results
or showing other applications of the tensorization techniques.
Though these topics were initially thought as `sidework',
some of them may be quite interesting for themselves.



\section{Appendix: On the norm of self-adjoint operators}%
\label{parOnTheNormOfSelfAdjointOperators}

In this appendix we prove the following
%
\begin{Lem}\label{l3253}
Let~$L$ be a self-adjoint operator on a real Hilbert space $H$, and let~$C<\infty$.
Then, to prove that $\VERT L \VERT \leq C$, it suffices to ensure that
\[\big\{ x\in H \colon \limsup_{k\longto\infty} |\langle L^k x, x \rangle|^{1/k} \leq C \big\}\]
is a dense subset of~$H$.
\end{Lem}

\begin{proof}
Reasoning by contraposition, we have to show that,
for~$L$ a self-adjoint operator on~$H$, for all~$C < \VERT L \VERT$,
the set of the~$x \in H$ such that
\[ \label{f8714} \limsup_{k\longto\infty} |\langle L^k x, x \rangle|^{1/k} > C \]
contains a non-empty open subset of~$H$.

Since $L$ is self-adjoint, by the spectral theorem~\cite[Theorem~7.18]{Weidmann},
it is unitarily equivalent to the ``multiplication by identity'' operator $M$
on a space $\bigoplus_{\alpha\in A} L^2(\rho_\alpha)$,
for~$A$ some set and~$\rho_\alpha$ some Radon measures on~$\RR$, that is
[in the following equation, the variable $\lambda$ is free, so that
$f(\lambda)$ is synonymous with~$f$]:
\[ M\Big( \sum_\alpha f_\alpha(\lambda) \Big) = \sum_\alpha \lambda f_\alpha(\lambda) .\]
So we will assume $L$ is of that form.

One has obviously:
\[\label{f7254} \VERT L \VERT = \sup \big\{ \lambda \geq 0 \colon 
(\exists \alpha \in A) \ \big(\rho_\alpha(\C{[-\lambda,\lambda]}) > 0\big) \big\} ;\]
moreover, for all~$f \in H$,
$f = \sum_{\alpha\in A} f_\alpha$ with $f_\alpha \in L^2(\rho_\alpha)$,
\[ \langle L^k f, f \rangle =
\sum_{\alpha\in A} \int_{\RR} \lambda^k |f_\alpha(\lambda)|^2 \,\dx\rho_\alpha(\lambda) ,\]
so that (observing that, for $k$ even, $\lambda^k \geq 0\ \ \forall\lambda$)
\[\label{f7360}
\limsup_{k\longto\infty} \big|\langle L^k f, f \rangle\big|^{1/k}
= \sup \Big\{\lambda \geq 0 \colon (\exists \alpha \in A) \ 
\Big(\int_{\C{[-\lambda,\lambda]}} |f_\alpha(\lambda')|^2 \,\dx\rho_\alpha(\lambda') > 0\Big)\Big\}.\]

Now, for $C < \VERT L \VERT$, the set
\[ U = \Big\{ f \in H \colon (\exists \alpha \in A) \ 
\Big( \int_{\C{[-C,C]}} |f_\alpha(\lambda)|^2 \,\dx\rho_\alpha(\lambda) > 0 \Big) \Big\} \]
is open because $\int_{\C{[-C,C]}} |f_\alpha(\lambda)|^2 \,\dx\rho_\alpha(\lambda)$
is a continuous function of~$f$, and it is non-empty by~(\ref{f7254}).
But (\ref{f8714}) is satisfied for all~$x\in U$ by~(\ref{f7360}),
so $U$ fulfills our quest.
\end{proof}

